% GENERATED FILE. Edit canonical lessons, then run npm run generate:books.
% canonical-source-hash: fnv1a64:b641262949bac0f6
% canonical-lessons: KA-C482-tore, KA-C482-dhumuku, KA-C482-more, KA-C482-cimmu, KA-C482-jinugu, KA-R482-review-476-482

\chapter{Streams and Waterfalls}
\label{ch:ka-streams-and-waterfalls}
\hlchaptermodality{482}

\begin{chapteropening}
\textbf{By the end of this chapter:} \emph{I can say a brook, to plunge, to rumble, to spurt and to trickle.}

Complete the last lesson of chapter 482: I can say a brook, to plunge, to rumble, to spurt and to trickle.
\end{chapteropening}

\section[\texorpdfstring{\kn{ತೊರೆ} (tore)}{ತೊರೆ (tore)}]{\kn{ತೊರೆ} (tore) — a brook; a stream}

\label{lesson:KA-C482-tore}



\begin{quote}
Before the new one: say the Kannada for a shell, then the Kannada for an estuary.
\end{quote}

\subsection*{You'll want to know: \kn{ತೊರೆ}}
\textbf{\kn{ತೊರೆ}} — \emph{tore} — \textquotedblleft{}a brook; a stream\textquotedblright{}.

\begin{quote}
\kn{ಕಾಡಿನ} \kn{ನಡುವೆ} \kn{ಒಂದು} \kn{ಚಿಕ್ಕ} \kn{ತೊರೆ} \kn{ಹರಿಯುತ್ತದೆ}.
\end{quote}

\begin{quote}
\emph{kāḍina naḍuve ondu cikka tore hariyuttade.}
\end{quote}

\begin{quote}
A small stream flows through the forest.
\end{quote}

Smaller than a \kn{ನದಿ}, like a \kn{ಝರಿ}.

\begin{cousinweb}
\textbf{\kn{ತೊರೆ}} (\emph{tore}) is native Dravidian: DEDR 3370 lists Kannada \emph{toṟe}, \textquotedblleft{}a stream, a river\textquotedblright{}.
\end{cousinweb}

\begin{grammarlens}[title={What you've built}]
One more word for getting around.
\end{grammarlens}

\subsection*{Your turn}
\begin{itemize}
\raggedright
  \item \emph{Say it:} \emph{tore}
  \item \emph{Say it:} \emph{tore}, once more
  \item \emph{Say it:} \emph{kāḍina naḍuve ondu cikka tore hariyuttade.} — A small stream flows through the forest.
  \item \emph{Recall:} say the Kannada for a shell, then the Kannada for an estuary, then say \emph{tore} again
\end{itemize}

\begin{tcolorbox}[breakable,colback=teal!4,colframe=teal!35!black,title={Before you move on}]
What does \kn{ತೊರೆ} mean? (\textquotedblleft{}A brook; a stream\textquotedblright{}.) Say it once more.
\end{tcolorbox}
\section[\texorpdfstring{\kn{ಧುಮುಕು} (dhumuku)}{ಧುಮುಕು (dhumuku)}]{\kn{ಧುಮುಕು} (dhumuku) — to plunge; to leap down}

\label{lesson:KA-C482-dhumuku}



\begin{quote}
Before the new one: say the Kannada for an estuary, then the Kannada for a brook.
\end{quote}

\subsection*{You'll want to know: \kn{ಧುಮುಕು}}
\textbf{\kn{ಧುಮುಕು}} — \emph{dhumuku} — \textquotedblleft{}to plunge; to leap down\textquotedblright{}.

\begin{quote}
\kn{ಜಲಪಾತದ} \kn{ನೀರು} \kn{ಬಂಡೆಯಿಂದ} \kn{ಕೆಳಗೆ} \kn{ಧುಮುಕುತ್ತದೆ}.
\end{quote}

\begin{quote}
\emph{jalapātada nīru baṇḍeyinda keḷage dhumukuttade.}
\end{quote}

\begin{quote}
The water of the falls plunges down from the rock.
\end{quote}

Past (I …): \kn{ಧುಮುಕಿದೆ}, \emph{dhumukide}.

Of water falling and of a diver: \kn{ನೀರಿಗೆ} \kn{ಧುಮುಕು}, to dive into water.

\begin{cousinweb}
\textbf{\kn{ಧುಮುಕು}} (\emph{dhumuku}) is native Dravidian: DEDR 3326 lists Kannada \emph{dumuku}, \textquotedblleft{}to leap or jump down from above\textquotedblright{}; the modern spelling writes it with the aspirate \kn{ಧ}.
\end{cousinweb}

\begin{grammarlens}[title={What you've built}]
One more word for getting around.
\end{grammarlens}

\subsection*{Your turn}
\begin{itemize}
\raggedright
  \item \emph{Say it:} \emph{dhumuku}
  \item \emph{Say it:} \emph{dhumuku}, once more
  \item \emph{Say it:} \emph{jalapātada nīru baṇḍeyinda keḷage dhumukuttade.} — The water of the falls plunges down from the rock.
  \item \emph{Recall:} say the Kannada for an estuary, then the Kannada for a brook, then say \emph{dhumuku} again
\end{itemize}

\begin{tcolorbox}[breakable,colback=teal!4,colframe=teal!35!black,title={Before you move on}]
What does \kn{ಧುಮುಕು} mean? (\textquotedblleft{}To plunge; to leap down\textquotedblright{}.) Say it once more.
\end{tcolorbox}
\section[\texorpdfstring{\kn{ಮೊರೆ} (more)}{ಮೊರೆ (more)}]{\kn{ಮೊರೆ} (more) — to rumble; to hum}

\label{lesson:KA-C482-more}



\begin{quote}
Before the new one: say the Kannada for a brook, then the Kannada for to plunge.
\end{quote}

\subsection*{You'll want to know: \kn{ಮೊರೆ}}
\textbf{\kn{ಮೊರೆ}} — \emph{more} — \textquotedblleft{}to rumble; to hum\textquotedblright{}.

\begin{quote}
\kn{ಮಳೆಗಾಲದಲ್ಲಿ} \kn{ಜಲಪಾತ} \kn{ಜೋರಾಗಿ} \kn{ಮೊರೆಯುತ್ತದೆ}.
\end{quote}

\begin{quote}
\emph{maḷegāladalli jalapāta jōrāgi moreyuttade.}
\end{quote}

\begin{quote}
In the rainy season the waterfall roars loudly.
\end{quote}

Past (it …): \kn{ಮೊರೆಯಿತು}, \emph{moreyitu}.

Of rushing water, bees and engines; a steady low sound.

\begin{cousinweb}
\textbf{\kn{ಮೊರೆ}} (\emph{more}) is native Dravidian: DEDR 4973 lists Kannada \emph{more}, \textquotedblleft{}to hum, to buzz, to murmur, to gurgle\textquotedblright{}, beside \emph{moraḷ}, \textquotedblleft{}to growl\textquotedblright{}.
\end{cousinweb}

\begin{grammarlens}[title={What you've built}]
One more word for getting around.
\end{grammarlens}

\subsection*{Your turn}
\begin{itemize}
\raggedright
  \item \emph{Say it:} \emph{more}
  \item \emph{Say it:} \emph{more}, once more
  \item \emph{Say it:} \emph{maḷegāladalli jalapāta jōrāgi moreyuttade.} — In the rainy season the waterfall roars loudly.
  \item \emph{Recall:} say the Kannada for a brook, then the Kannada for to plunge, then say \emph{more} again
\end{itemize}

\begin{tcolorbox}[breakable,colback=teal!4,colframe=teal!35!black,title={Before you move on}]
What does \kn{ಮೊರೆ} mean? (\textquotedblleft{}To rumble; to hum\textquotedblright{}.) Say it once more.
\end{tcolorbox}
\section[\texorpdfstring{\kn{ಚಿಮ್ಮು} (cimmu)}{ಚಿಮ್ಮು (cimmu)}]{\kn{ಚಿಮ್ಮು} (cimmu) — to spurt; to splash up}

\label{lesson:KA-C482-cimmu}



\begin{quote}
Before the new one: say the Kannada for to plunge, then the Kannada for to rumble.
\end{quote}

\subsection*{You'll want to know: \kn{ಚಿಮ್ಮು}}
\textbf{\kn{ಚಿಮ್ಮು}} — \emph{cimmu} — \textquotedblleft{}to spurt; to splash up\textquotedblright{}.

\begin{quote}
\kn{ಬಂಡೆಗೆ} \kn{ಬಡಿದ} \kn{ನೀರು} \kn{ಮೇಲಕ್ಕೆ} \kn{ಚಿಮ್ಮಿತು}.
\end{quote}

\begin{quote}
\emph{baṇḍege baḍida nīru mēlakke cimmitu.}
\end{quote}

\begin{quote}
Striking the rock, the water splashed up.
\end{quote}

Past (it …): \kn{ಚಿಮ್ಮಿತು}, \emph{cimmitu}.

Of water, blood or sparks that shoot out.

\begin{cousinweb}
\textbf{\kn{ಚಿಮ್ಮು}} (\emph{cimmu}) is native Dravidian: DEDR 2548 lists Kannada \emph{cimmu}, \textquotedblleft{}to flick, to squirt, to burst forth, to gush out\textquotedblright{}.
\end{cousinweb}

\begin{grammarlens}[title={What you've built}]
One more word for getting around.
\end{grammarlens}

\subsection*{Your turn}
\begin{itemize}
\raggedright
  \item \emph{Say it:} \emph{cimmu}
  \item \emph{Say it:} \emph{cimmu}, once more
  \item \emph{Say it:} \emph{baṇḍege baḍida nīru mēlakke cimmitu.} — Striking the rock, the water splashed up.
  \item \emph{Recall:} say the Kannada for to plunge, then the Kannada for to rumble, then say \emph{cimmu} again
\end{itemize}

\begin{tcolorbox}[breakable,colback=teal!4,colframe=teal!35!black,title={Before you move on}]
What does \kn{ಚಿಮ್ಮು} mean? (\textquotedblleft{}To spurt; to splash up\textquotedblright{}.) Say it once more.
\end{tcolorbox}
\section[\texorpdfstring{\kn{ಜಿನುಗು} (jinugu)}{ಜಿನುಗು (jinugu)}]{\kn{ಜಿನುಗು} (jinugu) — to trickle; to ooze}

\label{lesson:KA-C482-jinugu}



\begin{quote}
Before the new one: say the Kannada for to rumble, then the Kannada for to spurt.
\end{quote}

\subsection*{You'll want to know: \kn{ಜಿನುಗು}}
\textbf{\kn{ಜಿನುಗು}} — \emph{jinugu} — \textquotedblleft{}to trickle; to ooze\textquotedblright{}.

\begin{quote}
\kn{ಬಂಡೆಯ} \kn{ಮೇಲಿನಿಂದ} \kn{ನೀರು} \kn{ಜಿನುಗುತ್ತಿತ್ತು}.
\end{quote}

\begin{quote}
\emph{baṇḍeya mēlininda nīru jinuguttittu.}
\end{quote}

\begin{quote}
Water was trickling down from the top of the rock.
\end{quote}

Past (it …): \kn{ಜಿನುಗಿತು}, \emph{jinugitu}.

Of a thin drizzle too: \kn{ಮಳೆ} \kn{ಜಿನುಗುತ್ತಿದೆ}.

\begin{cousinweb}
\textbf{\kn{ಜಿನುಗು}} (\emph{jinugu}) is native Dravidian: DEDR 2520 lists Kannada \emph{jinugu}, with the noun \textquotedblleft{}drizzling rain\textquotedblright{}.
\end{cousinweb}

\begin{grammarlens}[title={What you've built}]
One more word for getting around.
\end{grammarlens}

\subsection*{Your turn}
\begin{itemize}
\raggedright
  \item \emph{Say it:} \emph{jinugu}
  \item \emph{Say it:} \emph{jinugu}, once more
  \item \emph{Say it:} \emph{baṇḍeya mēlininda nīru jinuguttittu.} — Water was trickling down from the top of the rock.
  \item \emph{Recall:} say the Kannada for to rumble, then the Kannada for to spurt, then say \emph{jinugu} again
\end{itemize}

\begin{tcolorbox}[breakable,colback=teal!4,colframe=teal!35!black,title={Before you move on}]
What does \kn{ಜಿನುಗು} mean? (\textquotedblleft{}To trickle; to ooze\textquotedblright{}.) Say it once more.
\end{tcolorbox}
\section[(dialogue)]{Looking back: from Wildlife Sanctuaries to Streams and Waterfalls}

\label{lesson:KA-R482-review-476-482}



\begin{quote}
Say the Kannada for wildlife and a sanctuary.
\end{quote}

\subsection*{Your turn}
Say these aloud:

\begin{itemize}
\raggedright
  \item wildlife, a sanctuary, a bird sanctuary, biodiversity, to migrate
  \item a wolf, a wild boar, to roar, to graze, a termite
  \item an eagle, a stork, a firefly, a crab, a whale
  \item vegetation, a bush, a grassland, to sprout, a swamp
  \item a rock, a desert, an open plain, a plateau, to wear away
  \item a coastline, a bay, coral, a shell, an estuary
  \item a brook, to plunge, to rumble, to spurt, to trickle
\end{itemize}

\begin{tcolorbox}[breakable,colback=teal!4,colframe=teal!35!black,title={Before you move on}]
Wildlife? (\textbf{\kn{ವನ್ಯಜೀವಿ}}, \emph{vanyajīvi}.) A sanctuary? (\textbf{\kn{ಅಭಯಾರಣ್ಯ}}, \emph{abhayāraṇya}.) A bird sanctuary? (\textbf{\kn{ಪಕ್ಷಿಧಾಮ}}, \emph{pakṣidhāma}.) Biodiversity? (\textbf{\kn{ಜೀವವೈವಿಧ್ಯ}}, \emph{jīvavaividhya}.) To migrate? (\textbf{\kn{ವಲಸೆ} \kn{ಹೋಗು}}, \emph{valase hōgu}.) A wolf? (\textbf{\kn{ತೋಳ}}, \emph{tōḷa}.) A wild boar? (\textbf{\kn{ಕಾಡುಹಂದಿ}}, \emph{kāḍuhandi}.) To roar? (\textbf{\kn{ಘರ್ಜಿಸು}}, \emph{gharjisu}.) To graze? (\textbf{\kn{ಮೇಯು}}, \emph{mēyu}.) A termite? (\textbf{\kn{ಗೆದ್ದಲು}}, \emph{geddalu}.) An eagle? (\textbf{\kn{ಹದ್ದು}}, \emph{haddu}.) A stork? (\textbf{\kn{ಕೊಕ್ಕರೆ}}, \emph{kokkare}.) A firefly? (\textbf{\kn{ಮಿಂಚುಹುಳು}}, \emph{miñcuhuḷu}.) A crab? (\textbf{\kn{ಏಡಿ}}, \emph{ēḍi}.) A whale? (\textbf{\kn{ತಿಮಿಂಗಿಲ}}, \emph{timiṅgila}.) Vegetation? (\textbf{\kn{ಸಸ್ಯ}}, \emph{sasya}.) A bush? (\textbf{\kn{ಪೊದೆ}}, \emph{pode}.) A grassland? (\textbf{\kn{ಹುಲ್ಲುಗಾವಲು}}, \emph{hullugāvalu}.) To sprout? (\textbf{\kn{ಚಿಗುರು}}, \emph{ciguru}.) A swamp? (\textbf{\kn{ಜೌಗು}}, \emph{jaugu}.) A rock? (\textbf{\kn{ಬಂಡೆ}}, \emph{baṇḍe}.) A desert? (\textbf{\kn{ಮರುಭೂಮಿ}}, \emph{marubhūmi}.) An open plain? (\textbf{\kn{ಬಯಲು}}, \emph{bayalu}.) A plateau? (\textbf{\kn{ಪ್ರಸ್ಥಭೂಮಿ}}, \emph{prasthabhūmi}.) To wear away? (\textbf{\kn{ಸವೆ}}, \emph{save}.) A coastline? (\textbf{\kn{ಕರಾವಳಿ}}, \emph{karāvaḷi}.) A bay? (\textbf{\kn{ಕೊಲ್ಲಿ}}, \emph{kolli}.) Coral? (\textbf{\kn{ಹವಳ}}, \emph{havaḷa}.) A shell? (\textbf{\kn{ಚಿಪ್ಪು}}, \emph{cippu}.) An estuary? (\textbf{\kn{ನದೀಮುಖ}}, \emph{nadīmukha}.) A brook? (\textbf{\kn{ತೊರೆ}}, \emph{tore}.) To plunge? (\textbf{\kn{ಧುಮುಕು}}, \emph{dhumuku}.) To rumble? (\textbf{\kn{ಮೊರೆ}}, \emph{more}.) To spurt? (\textbf{\kn{ಚಿಮ್ಮು}}, \emph{cimmu}.) To trickle? (\textbf{\kn{ಜಿನುಗು}}, \emph{jinugu}.)
\end{tcolorbox}
